\section{DECAF: A Semantic Decomposition of Paired Responses}
\label{sec:method}

Section~\ref{sec:setup} gives us two quantities: the response $r(t)$ at each
reveal level and the final contrast $d=r(1)$. \decaf uses the final contrast
only as a semantic reference for interpreting the intermediate responses.
If the final contrast is substantial, its direction tells us which responses
agree with the model's eventual behavior and which oppose it. If the final
contrast is negligible, intermediate responses cannot be interpreted as
supporting or opposing a meaningful final effect.

Choose a practical threshold $\epsilon>0$ and define
\[
a=\mathbf{1}\{|d|\ge \epsilon\},
\qquad
s=\operatorname{sign}(d).
\]
The gate $a$ determines whether the final contrast is active, while $s$
orients the trajectory when it is active. Let
\[
z(t)=s\,r(t),
\]
and write $z^+=\max\{z,0\}$ and $z^-=\max\{-z,0\}$.
DECAF routes the response as
\[
(e(t),c(t),f(t))
=
\bigl(a z^+(t),\,a z^-(t),\,(1-a)|r(t)|\bigr).
\]

An active response aligned with the final contrast becomes
\emph{evidence}; an active response pointing in the opposite direction
becomes \emph{contradiction}; and a response on a final-contrast-null pair
becomes \emph{fragility}. This logic is illustrated earlier in \figurename~\ref{fig:decaf-overview}.


Let $\mu$ be a normalized measure over stages and let the outer expectation cover paired examples and protocol randomness. We report
\begin{equation}
    M=\mathbb{E}|d|,
    \qquad
    (E,C,F)=\mathbb{E}\!\int_{\mathcal{T}}(e(t),c(t),f(t))\,\mathrm{d}\mu(t),
    \qquad
    \Abs=E+C+F.
    \label{eq:population-profile}
\end{equation}
The finite-grid implementation replaces the integral by quadrature. No additional model query is needed once the signed paired trajectory is available. See Appendices~\ref{app:theory} and~\ref{app:estimators} for finite-grid estimation, conditional summaries, threshold sensitivity, and
streaming accumulation.

\textbf{Forward-only implementation.}
The model enters only through evaluations of $q$. \decaf needs no gradients, parameters, or internal activations. It can run against a neural network, a tree ensemble, or a remote score API. Examples, stages, factors, paths, and models are independent batching dimensions.

\begin{figure}[H]
  \centering
  \begin{minipage}[t]{0.43\textwidth}
    \vspace{0pt}
    \centering
    \captionof{table}{Model calls per image in ImageNet-9. }
    \label{tab:complexity}
    \scriptsize
    \setlength{\tabcolsep}{2.5pt}
    \resizebox{\linewidth}{!}{%
\begin{tabular}{lrrrr}
\toprule
Method &
\shortstack{Fwd.} &
\shortstack{Bwd.} &
\shortstack{Internal} &
\shortstack{Learned} \\
\midrule

\rowcolor{decafblue}
\decaf-9
& 18
& \cellcolor{accessgreen}\textbf{0}
& \cellcolor{accessgreen}\textbf{No}
& \cellcolor{accessgreen}\textbf{No} \\

Input$\times$Grad
& 1 & 1 & Yes & No \\

IG
& 16 & 16 & Yes & No \\

SmoothGrad
& 16 & 16 & Yes & No \\

BlurIG
& 12 & 12 & Yes & No \\

Occlusion
& 49 & 0 & No & No \\

RISE
& 256 & 0 & No & No \\

\bottomrule
\end{tabular}
}
  \end{minipage}%
  \hfill
\begin{minipage}[t]{0.55\textwidth}
  \vspace{0pt}
  \centering
  \begin{tcolorbox}[
    enhanced,
    boxrule=0.5pt,
    colframe=algBorder,
    colback=algBG,
    arc=4pt,
    left=2pt, right=2pt, top=2pt, bottom=2pt,
    shadow={0mm}{0.3mm}{0mm}{gray!15} % 极轻微下落阴影
  ]
    \captionof{algorithm}{\textbf{DECAF in four steps.}}
    \label{alg:decaf}
    % \vspace{2pt}
    \small
    \begin{algorithmic}[1]
      \Require score $q$, paired path $(\mathbf{x}^{+}(t),\mathbf{x}^{-}(t))$, $\varepsilon$
      \State $d\gets q(\mathbf{x}^{+})-q(\mathbf{x}^{-})$; \;$a\gets\mathbf{1}\{|d|\geq\varepsilon\}$;\;$s\gets\operatorname{sign}(d)$
      \ForAll{$t\in\mathcal{T}$}\Comment{$r(t=1)=d$ is cached and reused}
        \State $r(t)\gets q(\mathbf{x}^{+}(t))-q(\mathbf{x}^{-}(t))$
        % 使用 colorbox 给核心三元组分解加上半透明背景高光
        \State \rlap{\colorbox{algHl}{\phantom{$(e,c,f)\gets\bigl(a(sr)^{+},a(sr)^{-},(1-a)|r|\bigr)$}}}$(e,c,f)\gets\bigl(a(sr)^{+},a(sr)^{-},(1-a)|r|\bigr)$
      \EndFor
      \State \Return $M,E,C,F$ by averaging over stages and examples
    \end{algorithmic}
  \end{tcolorbox}
\end{minipage}
\end{figure}

% Forward-only computation brings three practical advantages. Both branches and several stages can share a batch. Memory scales with inference rather than a backward graph. Large explanation jobs can be sharded across devices without synchronization. \emph{See Appendix~\ref{app:estimators} for query complexity, memory, batching, and black-box score-oracle implementations.}

\textbf{Scope of the semantics.}
\decaf assigns observable, endpoint-relative roles to paired responses; it does not infer the latent cause of a trajectory response. These response roles are defined relative to the score $q$, the chosen factual--counterfactual pair, the reveal trajectory, and the practical threshold  $\epsilon$. In particular,
\emph{fragility} denotes response on a pair whose final contrast is negligible
under this specification. It does not by itself distinguish off-manifold
artifacts, boundary uncertainty, calibration effects, or other possible causes.
An active pair may still be path-sensitive, but its response is described
through aligned and opposed mass rather than $F$.
